\chapter{Retrieval}
\label{ch:retrieval}

\section*{What this chapter is for}

\reported{Retrieval is the most frequently asked surface in this domain, and it
is the largest single category in analysed take-home
assessments.}{field-guide} It is also the most tutorialised, which means the
baseline answer is common and the differentiator is entirely in the failure
modes.

\section*{The questions}

\begin{bankq}{When would you \emph{not} use retrieval?}
\qasked
Frequently the opening question, and phrased as the negative on purpose.
\reported{Interviewers are reported to ask this form rather than asking what
retrieval is.}{field-guide-get-hired}

\qtested
Whether you reach for it by default. The question is a test of judgement
disguised as a test of knowledge, and answering it as though it were the
second is the failure.

\qstaff
Name the conditions under which it is the wrong tool, and give the alternative
for each. When the corpus is small enough to fit in the context window, put it
there \dash{} retrieval adds a failure mode and a latency budget to solve a
problem you do not have. When the question needs the whole corpus rather than
part of it \dash{} counting, aggregating, comparing across everything \dash{} no
top-$k$ will answer it, and the honest answer is a query against structured
data. When the knowledge is stable and small, it belongs in the instructions.
When the answer must be exact and auditable, a lookup beats a similarity
search. When the user's question and the document's wording do not share
vocabulary, plain semantic search will miss, and the fix is in the query or the
index rather than in more retrieval.

Then the summary that makes it one idea rather than a list: retrieval is for
when the relevant subset is small, identifiable from the question, and changes
often enough that you cannot bake it in.

\qsenior
A correct account of what retrieval is for \dash{} grounding, freshness,
citation, avoiding retraining \dash{} followed by one exception, usually the
context-window one. It is right. It answers the question that was not asked.

\qcontested
Where the boundary sits moves with context-window size and cost, and it has
been moving. \judgement{The reasonable position is that the boundary is now
economic rather than technical for a large class of corpora \dash{} you can
often put everything in, and whether you should is a cost and latency question.
Anybody who tells you the boundary is fixed is quoting a number from a
different year.}

\qdrill
Answer for a specific system: a product's documentation, twelve hundred pages,
updated weekly, questions from support staff. Then answer again with the corpus
at twelve pages. The answer should change completely, and if it does not, step
one of the method was skipped.
\end{bankq}

\begin{bankq}{How do you know your retrieval is any good?}
\qasked
Often the second question, and the one that separates candidates most sharply.

\qtested
Whether you have measured retrieval separately from generation. Most people
have not, and the give-away is an answer about output quality.

\qstaff
Separate the two. Retrieval quality is measurable on its own: build a set of
questions with the documents that should be found, and measure whether they
are found and where they rank. That number is independent of the model and does
not move when the prompt changes, which is what makes it useful \dash{} it is
the only part of the system you can regress against reliably.

Then the honest part: generation quality on top of good retrieval is a separate
measurement, and conflating them is why teams spend weeks tuning a prompt when
the document was never retrieved. \judgement{The first diagnostic question on
any bad answer is whether the right document was in the context at all, and
about half the time it was not.}

Building the question set is the work. It comes from real user questions where
you have them, and from documents read backwards where you do not \dash{} take
a passage, write the question it answers, and you have a labelled pair.

\qsenior
An answer about evaluating end-to-end output quality, possibly with a
model-based grader. It is not wrong and it measures the composition of two
systems, so a regression cannot be attributed to either.

\qcontested
Which retrieval metric to use is genuinely contested and mostly does not
matter. \judgement{Any rank-aware measure over a labelled set beats an
unlabelled opinion by a wide margin, and the difference between the candidate
metrics is much smaller than the difference between having a labelled set and
not.}

\qdrill
Describe how you would build the question set for a corpus you know, in ten
minutes of speech. Then say how many pairs you would want before believing the
number.
\end{bankq}

\begin{bankq}{Your answers are subtly wrong and the documents look right. What do you check?}
\qasked
A debugging question, usually posed as a scenario.

\qtested
Whether you have a diagnostic order or produce a list of everything that could
be wrong. The order is the answer.

\qstaff
Work from the outside in, because each step rules out a layer.

Was the right document retrieved at all? If not, it is an indexing or query
problem and nothing downstream matters. Was it retrieved but ranked below the
cut-off? That is a $k$ or a re-ranking problem. Was it in the context but
truncated at a chunk boundary, so the retrieved passage contains half the
answer? This is the commonest cause of \emph{subtly} wrong, and it is invisible
unless you look at the exact text that entered the context.

Was the whole answer present but contradicted by something else in the context
\dash{} an older document, a more specific one, a near-duplicate? Retrieval
returning several passages that disagree is a ranking problem wearing a
generation costume. And only when all of that is clean is it the model's
handling of correct context.

\judgement{The single most useful habit is logging the exact context as
assembled. Teams debug retrieval by looking at document identifiers and miss
that the passage was cut in half, because identifiers look right and text does
not lie.}

\qsenior
Adjusting chunk size, adding a re-ranker, or rewriting the instruction to say
``use only the provided context''. Each can work. Applied without the diagnosis
they are guesses, and the tell is that they are proposed in parallel rather
than in an order.

\qcontested
Chunking strategy is where most published advice lives and where the least
agreement exists. \judgement{Chunk-size advice that does not mention the shape
of the documents is worthless, and the only reliable method is to look at what
actually entered the context for the queries that failed.}

\qdrill
Take a system you have built or used and say, out loud, what you would need to
have logged to run this diagnosis. If the answer is ``more than I currently
log'', that is the finding.
\end{bankq}

\begin{bankq}{How would you cite sources, and what breaks?}
\qasked
Usually framed as a product requirement rather than a technical one.

\qtested
Whether you understand that a citation is a claim the system makes about
itself, and that it can be wrong independently of the answer.

\qstaff
The mechanism is straightforward: carry identifiers through with the passages
and require the output to reference them. What is interesting is the two
failure modes, both of which ship.

A citation can be present and unrelated \dash{} the model produces a plausible
answer and attaches a reference from the context that does not support it.
Nothing about the output looks wrong. This is only detectable by checking
whether the cited passage entails the sentence, which is itself an evaluation
problem.

And a citation can be correct while the answer is assembled from two passages
that should not have been combined. Each fragment is supported; the composition
is not.

\judgement{The reason this matters more than it appears is that citations are
sold to users as a trust mechanism. An unreliable citation is worse than none,
because it converts a system people would check into one they will not.}

\qsenior
The mechanism, correctly, plus prompt instructions to cite accurately. That
describes how to ask for citations rather than how to know you got them.

\qcontested
Whether citation accuracy should be checked automatically, and how.
\judgement{This is genuinely open. Automatic checking is itself model-based and
therefore has its own error rate, which puts you in a regress. The defensible
position is to sample rather than to gate, and to be honest with users about
what the citation means.}

\qdrill
Describe how you would measure citation accuracy on a system with no labels,
in under two minutes, naming what your method would miss.
\end{bankq}

\section*{What this chapter does not claim}

The frequency claim comes from one corpus with a stated method. It is the
strongest evidence available and it is one measurement.

The diagnostic order, the position on chunking, and the argument about
citations as a trust mechanism are the author's judgement, marked as such where
they appear. They are not consensus and the third in particular is a position
some teams would reject.
